\documentclass{article}
\usepackage[a4paper, margin=1in]{geometry}
\usepackage{graphicx} % Required for inserting images
\usepackage{xcolor}
\usepackage{hyperref}

\title{Research Questions for Studying LLM-Based POI Prediction}
\author{Lexi Liu}
\date{December 2025}

\begin{document}

\textbf{Reading list:}
\begin{enumerate}
    \item \href{https://web.stanford.edu/~jurafsky/slp3/}{Speech and Language Processing} by Dan Jurafsky and James H. Martin.\\
    Intuitive and easy to understand.
    
    \item Annotated version of the attention paper:  \href{https://nlp.seas.harvard.edu/annotated-transformer/}{https://nlp.seas.harvard.edu/annotated-transformer/}
    Helpful in coding attention-based models.

    \item UW CSE599 G1 \href{https://courses.cs.washington.edu/courses/cse493g1/24au/}{}
\end{enumerate}


Ranked from the easiest to the hardest:
\section{Semantic and Spatial Representation of POIs}


\begin{enumerate}
    \item How does the LLM represent POIs internally: as semantic entities or spatial locations?\\
    \textit{Background: POIs can be represented symbolically (using names or functions) or geometrically (using coordinates), 
    and LLMs may prefer one over the other. Experiments can ablate spatial information while preserving semantics to assess representational bias.}
    \begin{enumerate}
        \item How sensitive are predictions to the textual representation of POIs (e.g., POI IDs or functional categories)?\\
        \textit{Test robustness by varying POI naming schemes 
        (IDs, natural language names, categories) while holding trajectories fixed, 
        and measure prediction stability and accuracy.}
        \item Does classifying POIs into coarser semantic categories reduce prediction variance?\\
        \textit{Coarse categories may act as regularization ? compare variance and entropy of predictions under fine-grained versus aggregated POI labels.}
        
    \end{enumerate}

    
    \item How does the LLM interpret time in POI trajectories: 
    as a continuous process or 
    a discrete label?\\
    \textit{LLMs may treat time as symbolic tokens rather than continuous variables, affecting temporal generalization across resolutions.} \\
    Does providing numeric timestamps versus natural language time descriptions significantly change predictions or calibration?\\
    \textit{Evaluate performance and uncertainty under numeric times (e.g., Unix timestamps) versus descriptive times (e.g., ``late evening'').}
    
    \item To what extent does the LLM encode spatial proximity implicitly?\\
    
    \textit{Although LLMs lack explicit spatial embeddings, 
    co-occurrence patterns may induce implicit spatial structure. 
    This can be tested by removing all explicit distance cues? }
    
    \begin{enumerate}
        \item When given semantically similar but spatially distant candidate POIs, does the LLM favor semantic coherence or geographic proximity?\\
        \textit{Construct controlled choice sets where semantic similarity and spatial distance are decoupled, and analyze systematic selection preferences.}
        \item Does explicitly providing distance or travel-time information in the prompt help improve prediction accuracy?
    \end{enumerate}



    \item Does the LLM capture temporal dependence beyond immediate context?\\
    \textit{Long-range temporal dependencies are crucial for routine mobility; truncation experiments can reveal effective memory length.}\\
    How sensitive are predictions to addition of early trajectory history while keeping the recent window fixed (temporal dependence)?\\
    \textit{Incrementally add earlier visits and examine whether predictions shift, indicating non-Markovian reasoning.}

    \item How does the LLM handle periodic and cyclical mobility patterns?
    
    \textit{Many mobility processes exhibit strong diurnal or weekly cycles; compare LLM-generated visit frequencies against empirical seasonal patterns.}
    
    Does the LLM reproduce daily or weekly visitation cycles observed in the empirical data?

    \item {\color{blue}(Hard)} How does the LLM respond to structural changes in spatio-temporal dynamics?\\
    \textit{Simulate regime changes (e.g., lockdowns, relocation) and test whether the LLM adapts or continues extrapolating historical patterns.}
    Requires simulated or real regime shifts and nontrivial evaluation of adaptation behavior.

    \item {\color{blue}(Hard)}  Does the LLM jointly reason about space and time, or treat them independently?\\
    \textit{Joint reasoning would imply interaction effects between distance and time that cannot be reduced to marginal constraints.} Hard to evaluate or test.\\
    Does adding joint space--time constraints improve realism more than spatial or temporal constraints alone?\\
    \textit{Compare constraint violations and trajectory realism under joint versus marginal constraint prompting.}
\end{enumerate}

\section{Generative Behavior and Prediction}

\begin{enumerate}
    \item How often does the LLM select semantically plausible but empirically rare POIs under different prompting strategies?\\
    \textit{Quantify over-selection of low-frequency POIs and relate it to prompt specificity, examples, or likelihood-style prompting.}

    \item  {\color{blue}(Hard)} How do prediction errors accumulate in multi-step POI forecasting?\\
    \textit{Analyze error propagation?}
\end{enumerate}

\section{Prompt Design and Inductive Bias Control}

\begin{enumerate}
    \item How does the stated prediction objective influence LLM behavior?
    Variation across zero-shot, few-shot as well as fine-tuning?\\
    \textit{Different objectives may induce different inductive biases; compare accuracy, diversity, and constraint adherence across training regimes.}

    \item Can explicit spatial and temporal constraints improve prediction realism?\\
    \textit{Constraints may act as soft regularizers; evaluate improvements in physical plausibility without sacrificing predictive accuracy.}
    \begin{enumerate}
        \item How often does the LLM violate explicit distance or time constraints under constrained prompts?\\
        \textit{Measure violation rates and severity to assess whether constraints are internally enforced or ignored.}
    \end{enumerate}


    \item ({\color{blue}Hard}) Can probabilistic prompting elicit calibrated uncertainty estimates?\\
    \textit{Probabilistic prompts aim to align LLM outputs with predictive uncertainty rather than point predictions.}
    \begin{enumerate}
        \item How well do elicited probabilities align with empirical frequencies across POI categories?\\
        \textit{Evaluate calibration using reliability diagrams or expected calibration error across frequency strata.}
        \item Does predicted uncertainty increase appropriately in ambiguous or low-frequency contexts?\\
        \textit{Test whether entropy or variance increases for sparse histories or multiple plausible next locations.}
    \end{enumerate}
\end{enumerate}

\section{Failure Modes and Stress Tests}

\begin{enumerate}
    \item How does the LLM perform on rare or atypical mobility behaviors?\\
    \textit{Stress-test the model on outliers to identify regression toward typical routines.}
    \begin{enumerate}
        \item Are irregular schedules or night-time trajectories systematically mapped to typical routines?\\
        \textit{Compare prediction diversity and timing bias for atypical versus regular users.}
        \item Does the LLM hallucinate common POIs in sparsely observed regimes?\\
        \textit{Track hallucination rates when historical data are deliberately thinned or partially masked.}
    \end{enumerate}

    \item 
    How does performance change for unseen POIs or newly developed regions?\\
    \textit{Evaluate zero-shot generalization by holding out regions or POIs and measuring semantic versus spatial extrapolation.}
\end{enumerate}

\section{Experiment}


\noindent 
The key to a dynamic GCN is to learn a weight matrix that changes through time.
Consider a GCN equation:
at time $t$,
the $\ell$th layer 
update a node embedding matrix 
by
\[
H^{(\ell+1)}_t = \sigma(P_t H^{(\ell)}_t W^{(\ell)}_t) 
\]
where $P_t$ is a normalized adjacency matrix defined as 
$P_t = D_t^{-1/2} A_t D_t^{-1/2}$
where $D_t$ is the degree matrix 
defined as 
\[
D^{(i,j)}_t = 
\begin{cases}
    \deg (v^{(i)}_t) 
    \quad \text{if}\quad i=j;\\
    0 \quad \text{otherwise}
\end{cases}
\]
where the degree $\deg (v^{(i)}_t)$ of a vertex counts the number of times an edge terminates at that vertex. 
The function $\sigma(\cdot)$
is an activation function 
for all layers except for the output layer.
An example is the ReLU function $\sigma(x) = \max(x,0)$.
The node embedding matrix is initialized by the input node features,
i.e.,
$H^{(0)}_t = X_t$
where $X_t \in \mathbb{R}^{v_t}$
is a vector containing the visitor counts at time $t$.\\



\subsection{Learned graph structure with MTGNN}

Let $E_1, E_2 \in \mathbb{R}^{N \times d}$ denote learnable node embeddings. 
The adjacency matrix is defined as
\[
A = \mathrm{softmax}\big(\mathrm{ReLU}(E_1 E_2^\top)\big).
\]
This formulation induces a directed, weighted graph, where edge weights reflect similarity in a latent embedding space. The row-wise softmax normalization ensures numerical stability and facilitates interpretation as transition weights. This component enables the model to uncover hidden relationships among variables when the underlying graph is unknown.\\

\noindent 
For modeling the temporal structure,
the model uses dilated inception CNN,
which is designed to capture long-range temporal dependencies.
The algorithm is summarized as follows:
\begin{algorithm}[H]
\caption{MTGNN for Multivariate Time Series Forecasting \citep{wu2020connecting}}
\label{alg:mtgnn}
\begin{algorithmic}[1]
\REQUIRE Input time series $X \in \mathbb{R}^{T \times N}$, 
input length $L$, prediction horizon $H$, 
number of layers $K$

\ENSURE Predictions $\hat{Y} \in \mathbb{R}^{H \times N}$

\STATE Extract sliding window input:
\[
X_{\text{in}} = X_{(t-L+1):t}
\]

\STATE Learn adaptive adjacency matrix:
\[
A = \text{softmax}\big(\text{ReLU}(E_1 E_2^\top)\big)
\]
where $E_1, E_2 \in \mathbb{R}^{N \times d}$ are learnable node embeddings

\STATE Initialize hidden representation:
\[
H^{(0)} = X_{\text{in}}
\]

\FOR{$k = 1$ to $K$}

    \STATE Temporal convolution (dilated causal CNN):
    \[
    \tilde{H}^{(k)} = \text{TCN}\big(H^{(k-1)}\big)
    \]

    \STATE Graph convolution (mix-hop propagation):
    \[
    H^{(k)} = \sum_{m=0}^{M} A^m \tilde{H}^{(k)} W_m^{(k)}
    \]

\ENDFOR

\STATE \textbf{Output layer:}
\[
\hat{Y} = \text{Linear}(H^{(K)})
\]

\RETURN $\hat{Y}$
\end{algorithmic}
\end{algorithm}

\noindent The key component for accomodating the evolving graphs is the update of the
weight matrix $W^{(\ell)}_t$ at time $t$ based on the
current as well as
historical information. 
This requirement can be naturally fulfilled by using a recurrent architecture.
One option is to use the  long short-term memory (LSTM)
cell to model this input-output relationship
written as 
\[
W^{(\ell)}_t = \text{LSTM}(W^{(\ell)}_{t-1}).
\]
The LSTM structure is used
to learn how the convolution rule should evolve.
Our prediction is generated after an output layer
\begin{equation}\label{eqt:output}
    \hat{X}_{T+1} = H_T W_{\textrm{out}} 
\end{equation}
where $W_{\textrm{out}} \in \mathbb{R}^{v_{T}}$ is a 
weight matrix.
The algorithm is described as follows:

\begin{algorithm}[H]
\caption{EvolveGCN-H for Dynamic Graph Learning \citep{pareja2020evolvegcn}}
\label{alg:evolvegcnH}
\begin{algorithmic}[1]
\REQUIRE Sequence of graphs $\{G_t = (A_t, X_t)\}_{t=1}^T$, 
initial GCN weights $W_0^{(l)}$ for layers $l = 1, \dots, L$
\ENSURE Node embeddings $\{H_t^{(l)}\}$

\FOR{$t = 1$ to $T$}
    \STATE Initialize $H_t^{(0)} = X_t$
    
    \FOR{$l = 1$ to $L$}
        \STATE \textbf{Evolve weights via GRU:}
        \[
        W_t^{(l)} = \text{GRU}\big(H_{t-1}^{(l)}, \, W_{t-1}^{(l)}\big);
        \]
        
        \STATE \textbf{Graph convolution:}
        \[
        H_t^{(l)} = \sigma\big(\tilde{A}_t H_t^{(l-1)} W_t^{(l)}\big),
        \]
        where $\tilde{A}_t$ denotes the normalized adjacency matrix $A_t$ at time $t$.
    \ENDFOR
\ENDFOR

\RETURN Output $\hat{X}_{T+1}$
using (\ref{eqt:output}).
\end{algorithmic}
\end{algorithm}


\subsection{Pre-determined adjacency matrices}

\noindent
To use the EvolveGCN framework in \cite{pareja2020evolvegcn},
we need to determine the adjacency matrices $A_t$ as requirements.
The ways we define the adjacency matrix depend on what kind of relationship between the POIs we would like to capture.
We explore several options
for determining the adjacency matrix $A_t$ at time $t$:
\begin{itemize}
    \item(Distance-based graph) We can define adjacency via geographic distance as 
    \[
        A_{ij} = \exp\left(-\frac{d(i,j)^2}{\sigma^2}\right)
    \]
    where $d(i,j)$ is a distance measure between two POIs.
    We use the Haversine distance to compute the distance measuer, i.e.,
    \[
    d(i,j) = 2r\arcsin\sqrt{\sin^2(\Delta \phi/2) + \cos(\phi_i)\cos(\phi_j)\sin^2(\Delta \lambda/2)}
    \]
    where $\Delta \phi = (\phi_i - \phi_j)\times \pi/180$,
    $\Delta \lambda = (\lambda_i - \lambda_j)\times \pi/180$
    are the differences of latitudes represented by $\phi_i, \phi_j$
    and of longitudes $\lambda_i, \lambda_j$.
    The radius of earth $r=6371$km approximately.

    \item(Semantic similarity graph)
    The adjacency matrix can also be defined to describe the semantic similarity of the POIs.
    For example, we can consider
    \[
    A_{ij} = \mathds{1}(\text{category}_i=\text{category}_j)
    \]
    where $\text{category}_i$
    is a label of the $i$th POI.

\end{itemize}


\end{document}
